\section{Position norms and a controlled group average}
\label{sec:orbit}
Let $X=\mathrm{dS}_2$ with global coordinates $(t,\theta)$ and
\begin{equation}
 \dd x=\cosh t\,\dd t\,\dd\theta,\qquad \theta\in[0,2\pi).
\end{equation}
Let $\mathcal K$ be a separable Hilbert space carrying a strongly continuous
unitary representation $V$ of $G$. It will subsequently be a fixed
matter-particle sector. On $Y=X\times X$ the diagonal action is
$(U(g)\psi)(y)=V(g)\psi(g^{-1}y)$.

Outside the null set where the two embedding vectors have inner product
$\pm1$, the position pairs form six free orbit families. Take the left point
to be $(0,0)$ and the right point to be one of
\begin{equation}
 (\pm\tau,0),\ (\pm\tau,\pi),\quad \tau>0;
 \qquad (0,\pm\alpha),\quad 0<\alpha<\pi.
\end{equation}
Writing this disjoint set of sections as $Z$ and its section map as $s(z)$,
the measure decomposition is
\begin{equation}
 \dd y=\dd g\,\dd\mu(z),\qquad
 \dd\mu=\sinh\tau\,\dd\tau\quad\text{or}\quad\sin\alpha\,\dd\alpha.
 \label{eq:coarea}
\end{equation}
The normalization agrees with $\dd g=\dd\theta_1\dd\theta_2
\sinh\lambda\dd\lambda$. Near the identity the two-boost/rotation
Jacobian for the first type of section has absolute value $\tanh\tau$;
multiplication by the position density $\cosh\tau$ gives $\sinh\tau$.
The spacelike section gives $|\sin\alpha|$. This fixes the constant
without dividing by an infinite group volume.

Define the ambient gauge-fixed space $\cF(\mathcal K)=L^2(Z,\dd\mu;\mathcal K)$.
For $\psi\in L^2(Y;\mathcal K)$ set
\begin{align}
 F_\psi(g,z)&=V(g^{-1})\psi(gs(z)),\\
 a_\psi(z)&=\int_G\norm{F_\psi(g,z)}\dd g,\\
 \cD_{\mathrm{orb}}&=\{\psi\in L^2(Y;\mathcal K):a_\psi\in L^2(Z)\}.
\end{align}
On this domain define $B\psi(z)=\int_GF_\psi(g,z)\dd g$ as a Bochner
integral and its invariant extension by $(\eta\psi)(gs(z))=V(g)B\psi(z)$.

\begin{proposition}[Orbit Fubini identity]
For $\phi,\psi\in\cD_{\mathrm{orb}}$, the group average is absolutely
convergent and
\begin{equation}
 \int_G\ip{\phi}{U(g)\psi}\dd g=\ip{B\phi}{B\psi}_{\cF(\mathcal K)}.
 \label{eq:fp-ga}
\end{equation}
For inverse-symmetric compact cutoffs $0\le\chi_R\le1$ tending to one,
the orbit integrals $B_R\psi$ converge to $B\psi$ in $\cF(\mathcal K)$.
\end{proposition}
\begin{proof}
Unimodularity and a Haar change of variables give
\begin{equation}
 \int_G\bigl|\ip{\phi}{U(g)\psi}\bigr|\dd g
 \le\int_Za_\phi(z)a_\psi(z)\dd\mu(z)
 \le\norm{a_\phi}_2\norm{a_\psi}_2.
\end{equation}
Fubini separates the two Bochner integrals and gives \eqref{eq:fp-ga}.
The cutoff error is bounded by
$\norm{\int_G(1-\chi_R(g))\norm{F_\psi(g,\cdot)}\dd g}_2$,
which tends to zero by dominated convergence. A finite-cutoff group average
is generally not invariant. An invariant extension of its section values is
a different object and is not presumed to retain positive frequency.
\end{proof}

For single angular modes $u_L,u_R$ and $v\in\mathcal K$, let
$f_i=|u_i|$ and $A_{f_i}(g)=\ip{f_i}{U_X(g)f_i}$.
Nonnegative Tonelli and \eqref{eq:coarea} give the exact identity
\begin{equation}
 \norm{a_{u_Lu_Rv}}_2^2
 =\norm v^2\int_G A_{f_L}(g)A_{f_R}(g)\dd g.
 \label{eq:orbit-autocorrelation}
\end{equation}
No decay assumption on a matter matrix coefficient is used.

\begin{lemma}[Finite position states are distributions]
For every $k\in\cF(\mathcal K)$, the extension
$Rk(gs(z))=V(g)k(z)$ lies in $L^1_{\mathrm{loc}}(Y;\mathcal K)$.
The map $R$ is continuous and injective. For every $\phi\in\cD_{\mathrm{orb}}$,
\begin{equation}
 T_k(\phi):=\int_Y\ip{Rk(y)}{\phi(y)}\dd y
 =\ip{k}{B\phi},\qquad |T_k(\phi)|\le\norm k_2\norm{a_\phi}_2.
 \label{eq:position-pairing}
\end{equation}
\label{lem:position-distribution}
\end{lemma}
\begin{proof}
The radial bound below implies that a compact smooth position test belongs
to $\cD_{\mathrm{orb}}$, with $\norm{a_\phi}_2\le C_C\sup_C\norm\phi$
on a fixed compact set $C$: dominate its norm by two compact radial
envelopes and apply \eqref{eq:orbit-autocorrelation}.
Tonelli gives
\begin{equation}
 \int_Y\norm{Rk(y)}\norm{\phi(y)}\dd y
 =\int_Z\norm{k(z)}a_\phi(z)\dd\mu(z).
\end{equation}
A compact cutoff equal to one on $C$ proves the local $L^1$ bound and
\eqref{eq:position-pairing}. If $Rk=0$ almost everywhere, integrate over
the image of a finite positive Haar-measure set times $Z$ in
\eqref{eq:coarea}; then $k=0$ almost everywhere. No trace of an arbitrary
$L^1$ function on a section is taken.
\end{proof}

For a smooth noncompact orbit test, \eqref{eq:position-pairing} is also the
limit of compact smooth position cutoffs. The orbit norm of
$(1-\xi_j)\phi$ tends to zero when $0\le\xi_j\le1$ tend to one.
Thus these pairings are determined by the position distribution.